\subsection{Group comparisons}\label{sec:val-effects}

\paragraph{The effect likelihood.} The engine with age segments (Section~\ref{sec:likelihood}) reproduces the
calibrator's own likelihood on 26 control cohorts of the NIA Interventions Testing Program to
\C{eff-seg-vs-j-per-ind} nats per individual, and an independent piecewise implementation to
\C{eff-seg-vs-srk-per-ind} nats per individual, for changes of $\eta$, $\beta$, $\eps$, $\kappa$, $A$, $\xc$ and $h$
in arms that include a treatment started and stopped in mid-life. Its exact gradients agree with finite differences to
a relative \C{eff-grad-rel}.
% src: aging:code_5/itp_analysis/recal/results/gates/V1_j100.json, V2.json, V3.json
A threshold change that starts at birth can be computed either with a threshold segment or by the rescaling of
Section~\ref{sec:model}; the two agree to \C{eff-gx-birth-max} nats per individual. For a change that starts at
$t_0>0$ they differ by \C{eff-gx-late-min}--\C{eff-gx-late-max} nats per individual, because the rescaling also
rescales the damage accumulated before $t_0$. Late threshold changes are therefore always computed with the segment.
% src: aging:code_4/Mire_fly_interventions/recal_2026-10/results/p3/gates/gates.json (G-X)

\paragraph{Recovery with independent individuals.} A pre-registered recovery study simulated arms and controls at the
designs of the Interventions Testing Program, with independent animals, from a null and from single-parameter and
two-parameter truths, and applied the procedure of Section~\ref{sec:effects} with the reference fixed at its fit.
\begin{itemize}[leftmargin=1.5em,itemsep=1pt]
  \item Nominal 95\,\% profile intervals for the multiplier covered the truth in a fraction \C{itp-coverage} of
    \C{itp-coverage-n} intervals.
  \item At the asymptotic Šidák threshold, the null was rejected in a fraction \C{itp-fpr-8arm}--\C{itp-fpr-1arm} of
    replicates, depending on the number of arms pooled, i.e. at or below the nominal 0.05.
  \item A single-parameter truth was identified correctly in a fraction \C{itp-ident-single-1} of replicates with one
    arm, rising to \C{itp-ident-single-2}, \C{itp-ident-single-4} and \C{itp-ident-single-8} with two, four and eight
    arms that share the change. Two-parameter truths were identified in only \C{itp-ident-pair-1}--\C{itp-ident-pair-8}
    when each arm had its own effect size.
  \item Naming the group instead of the parameter (production $\{\eta,A\}$, removal $\{\beta,\kappa\}$, noise and
    threshold $\{\eps,\xc\}$; Section~\ref{sec:effects}) raised the share of correct answers for single truths from
    \C{grp-weak1-exact} to \C{grp-weak1-group} with one arm at a weak effect (expected $2\Delta\ell=4$), from
    \C{grp-weak8-exact} to \C{grp-weak8-group} with eight arms, and from \C{grp-obs1-exact} to \C{grp-obs1-group} with
    one arm at an effect that moves the median lifespan by 15\,\%. Most of the gain comes from $\eta$ confused with $A$
    and $\kappa$ with $\beta$. Removal is the weakest group: with eight arms at the weak size its right-group rate was
    \C{grp-removal-weak8-lo}--\C{grp-removal-weak8-hi}, against at least \C{grp-other-weak8-lo} for the other
    parameters.
\end{itemize}
% src: aging:code_5/itp_analysis/recal/results/2026-10-18_J100/pool_pairs/ (pool_groups_singles.csv; Stage 5 replicates)

\paragraph{Pooling arms to identify pairs.} A second simulation study asked how many pooled arms, and of what kind, make
a two-parameter change identifiable. Two ITP arms served as bases, each with its own control, which was refitted in
every replicate. Five pair truths were simulated, and eight arms were pooled
in three designs (each pair member moving the median by about 7.5\,\%): arms spread over four dosing ages with one shared effect size, the same with a separate size per arm,
and eight repeats at one dosing age with separate sizes (\C{pool-reps} replicates per cell, exact fits).
\begin{itemize}[leftmargin=1.5em,itemsep=1pt]
  \item Spreading the arms over dosing ages with one shared size was the best design. It named $A+\eps$ in
    \C{pool-aeps-lo}--\C{pool-aeps-hi} of replicates (at most \C{pool-aeps-repeats-hi} with repeats), $A+\xc$ in
    \C{pool-axc-lo}--\C{pool-axc-hi}, and $\eta+A$ in \C{pool-etaa-lo}--\C{pool-etaa-hi}.
  \item Production against removal did not separate: $\eta+\beta$ was named in at most \C{pool-etabeta-max} of
    replicates under any design, about the rate of a random choice among the candidate pairs.
  \item At the level of groups, the best design named the set of changed groups in \C{grp-pairs-best-lo}--\C{grp-pairs-best-hi}
    of replicates, and at least one changed group in at least \C{grp-pairs-onegroup-min} for every truth.
\end{itemize}
% src: aging:code_5/itp_analysis/recal/results/2026-10-18_J100/pool_pairs/ (README.md, pool_sim_identification.csv, pool_groups_pairs.csv)
A pair can thus be named from pooled data, but only from a design built for it, and production against removal could
not be separated at eight arms.
% src: aging:code_5/itp_analysis/recal/docs/PREREG_RECOVERY.md and recal/results/2026-10-18_J100/recovery_verdict/
A single arm of a typical mouse lifespan study can therefore usually tell \emph{that} survival changed, but often not
\emph{which} parameter changed. This is why condition 6 of the naming rule asks for recall and precision at the
study's own design.

\paragraph{The analytic shortcut.} The expected log-likelihood gap $D_{pq}$ (Section~\ref{sec:effects}) was checked
against recovery simulations on \C{eff-ia-cells} design cells, totalling \C{interv-replicates} simulated replicates.
$D$ and the simulated confusion rate were correlated as expected (Spearman $\rho=\C{eff-ia-rho}$). In all
\C{eff-ia-cells-d9} cells with $D\ge9$ the confusion rate was at most \C{eff-ia-confusion-max}
(\C{eff-ia-violations} violations).
% src: aging:interventions_data/calib/ident/ia_validation/summary.json

\paragraph{Clustered individuals.} One pre-registered simulation study used clustered individuals: a fly
dietary-restriction design with flies housed in cages and cage effects at the measured between-cage variance.
% src: aging:code_4/Mire_fly_interventions/recal_2026-10/docs/PREREG_P3.md, results/p3/P3_SUMMARY.md
\begin{itemize}[leftmargin=1.5em,itemsep=1pt]
  \item Dividing the likelihood-ratio statistic by a Pearson dispersion $\phi$ and comparing it with $\chi^2$ rejected
    the null in \C{fly-null-hits} of \C{fly-null-n} null replicates, a fraction \C{fly-null-rate} (Wilson 95\,\%
    interval \C{fly-null-wilson-lo}--\C{fly-null-wilson-hi}), against a nominal 0.05.
  \item The spread of the estimated multipliers across replicates was larger than the $\phi$-widened intervals imply:
    those nominal 95\,\% intervals would cover the truth in only about \C{fly-cov-441}--\C{fly-cov-707} of
    replicates. The bias of the estimates was at most \C{fly-bias} in $\log_{10}$.
  \item With the null threshold calibrated by simulation instead, the one truth simulated in full (a production
    change) was recovered with recall \C{fly-recall}; precision could not be assessed, since no other truth was
    simulated at that design.
\end{itemize}
Design-effect scaling is therefore not a substitute for a calibrated threshold when individuals share cages, vials or plates, which
is the basis of the recommendation in Section~\ref{sec:effects}.


